\documentclass[11pt]{article}

\usepackage[letterpaper,margin=0.82in,headheight=15pt]{geometry}
\usepackage[T1]{fontenc}
\usepackage{lmodern}
\usepackage{microtype}
\usepackage{amsmath,amssymb,amsthm,bm,mathtools}
\usepackage{stmaryrd}
\usepackage{booktabs,longtable,tabularx,array,multirow}
\usepackage{enumitem}
\usepackage{xcolor}
\usepackage{fancyhdr}
\usepackage{graphicx}
\usepackage{tikz-cd}
\usepackage{hyperref}
\usepackage[nameinlink,noabbrev]{cleveref}

\definecolor{navy}{HTML}{17324D}
\definecolor{teal}{HTML}{147D92}
\definecolor{orange}{HTML}{D47A27}
\definecolor{softblue}{HTML}{EEF4F7}
\definecolor{softgray}{HTML}{F3F5F6}
\definecolor{darkgray}{HTML}{46545E}
\definecolor{softorange}{HTML}{FFF4E8}
\definecolor{softgreen}{HTML}{EDF7F0}
\definecolor{softred}{HTML}{FCEEEE}

\hypersetup{
  colorlinks=true,
  linkcolor=navy,
  citecolor=teal,
  urlcolor=teal,
  pdftitle={Volume XIII: Spin-Resolved NNpi Dynamics, Pion Matching, Two-Body Currents, and Coherent Nuclear Amplitudes},
  pdfauthor={DeuteronWigner project}
}

\pagestyle{fancy}
\fancyhf{}
\fancyhead[L]{\small\color{darkgray}DeuteronWigner formalism - Volume XIII}
\fancyhead[R]{\small\color{darkgray}Spin-resolved NNpi and coherent nuclear dynamics}
\fancyfoot[L]{\small\color{darkgray}31 July 2026}
\fancyfoot[R]{\small\color{darkgray}\thepage}
\renewcommand{\headrulewidth}{0.4pt}

\setlist[itemize]{leftmargin=1.5em,itemsep=2pt,topsep=3pt}
\setlist[enumerate]{leftmargin=1.7em,itemsep=2pt,topsep=3pt}
\setlength{\parindent}{0pt}
\setlength{\parskip}{5pt}
\setlength{\emergencystretch}{2em}
\renewcommand{\arraystretch}{1.16}
\allowdisplaybreaks

\newtheorem{definition}{Definition}[section]
\newtheorem{axiom}[definition]{Axiom}
\newtheorem{principle}[definition]{Design principle}
\newtheorem{requirement}[definition]{Requirement}
\newtheorem{proposition}[definition]{Proposition}
\newtheorem{theorem}[definition]{Theorem}
\newtheorem{lemma}[definition]{Lemma}
\newtheorem{corollary}[definition]{Corollary}
\newtheorem{remark}[definition]{Remark}
\newtheorem{decision}[definition]{Model decision}

\crefname{definition}{definition}{definitions}
\crefname{axiom}{axiom}{axioms}
\crefname{principle}{design principle}{design principles}
\crefname{requirement}{requirement}{requirements}
\crefname{proposition}{proposition}{propositions}
\crefname{theorem}{theorem}{theorems}
\crefname{lemma}{lemma}{lemmas}
\crefname{corollary}{corollary}{corollaries}
\crefname{remark}{remark}{remarks}
\crefname{decision}{model decision}{model decisions}

\newcommand{\kT}{\bm{k}_{T}}
\newcommand{\pT}{\bm{p}_{T}}
\newcommand{\qT}{\bm{q}_{T}}
\newcommand{\lT}{\bm{\ell}_{T}}
\newcommand{\rT}{\bm{r}_{T}}
\newcommand{\sT}{\bm{s}_{T}}
\newcommand{\DT}{\bm{\Delta}_{T}}
\newcommand{\bD}{\bm{b}_{\Delta}}
\newcommand{\bM}{\bm{b}_{\mathrm{TMD}}}
\newcommand{\RT}{\bm{R}_{T}}
\newcommand{\dd}{\mathrm{d}}
\newcommand{\Tr}{\operatorname{Tr}}
\newcommand{\tr}{\operatorname{tr}}
\newcommand{\diag}{\operatorname{diag}}
\newcommand{\rank}{\operatorname{rank}}
\newcommand{\Span}{\operatorname{span}}
\newcommand{\Ker}{\operatorname{Ker}}
\newcommand{\Ran}{\operatorname{Ran}}
\newcommand{\Cov}{\operatorname{Cov}}
\newcommand{\Var}{\operatorname{Var}}
\newcommand{\Id}{\operatorname{Id}}
\newcommand{\Hom}{\operatorname{Hom}}
\newcommand{\End}{\operatorname{End}}
\newcommand{\Hil}{\mathcal H}
\newcommand{\LF}{\ensuremath{\mathrm{LF}}}
\newcommand{\TMD}{\ensuremath{\mathrm{TMD}}}
\newcommand{\GTMD}{\ensuremath{\mathrm{GTMD}}}
\newcommand{\GPD}{\ensuremath{\mathrm{GPD}}}
\newcommand{\PDF}{\ensuremath{\mathrm{PDF}}}
\newcommand{\QCD}{\ensuremath{\mathrm{QCD}}}
\newcommand{\EFT}{\ensuremath{\mathrm{EFT}}}
\newcommand{\TTN}{\ensuremath{\mathrm{TTN}}}
\newcommand{\TN}{\ensuremath{\mathrm{TN}}}
\newcommand{\PCAC}{\ensuremath{\mathrm{PCAC}}}
\newcommand{\cO}{\mathcal O}
\newcommand{\cM}{\mathcal M}
\newcommand{\cR}{\mathcal R}
\newcommand{\cS}{\mathcal S}
\newcommand{\cT}{\mathcal T}
\newcommand{\cV}{\mathcal V}
\newcommand{\cW}{\mathcal W}
\newcommand{\cE}{\mathcal E}
\newcommand{\cG}{\mathcal G}
\newcommand{\cH}{\mathcal H}
\newcommand{\cK}{\mathcal K}
\newcommand{\cQ}{\mathcal Q}
\newcommand{\cP}{\mathcal P}
\newcommand{\cC}{\mathcal C}
\newcommand{\cZ}{\mathcal Z}
\newcommand{\cD}{\mathcal D}
\newcommand{\cN}{\mathcal N}
\newcommand{\cU}{\mathcal U}
\newcommand{\cF}{\mathcal F}
\newcommand{\cA}{\mathcal A}
\newcommand{\cB}{\mathcal B}
\newcommand{\cL}{\mathcal L}
\newcommand{\cJ}{\mathcal J}
\newcommand{\cI}{\mathcal I}
\newcommand{\eps}{\varepsilon}
\newcommand{\abs}[1]{\left|#1\right|}
\newcommand{\norm}[1]{\left\lVert#1\right\rVert}
\newcommand{\inner}[2]{\left\langle #1\middle|#2\right\rangle}
\newcommand{\avg}[1]{\left\langle #1\right\rangle}
\newcommand{\phys}{\mathrm{phys}}
\newcommand{\eff}{\mathrm{eff}}
\newcommand{\bare}{\mathrm{bare}}
\newcommand{\ind}{\mathrm{ind}}
\newcommand{\ct}{\mathrm{ct}}
\newcommand{\reg}{\mathrm{reg}}
\newcommand{\trunc}{\mathrm{trunc}}
\newcommand{\ren}{\mathrm{ren}}
\newcommand{\num}{\mathrm{num}}
\newcommand{\UV}{\mathrm{UV}}
\newcommand{\IR}{\mathrm{IR}}
\newcommand{\COM}{\mathrm{CM}}
\newcommand{\Lz}{L_z}
\newcommand{\Jz}{J^z}
\newcommand{\Amp}{\mathbf{Amp}}
\newcommand{\Dens}{\mathbf{Dens}}
\newcommand{\Match}{\mathbf{Match}}
\newcommand{\Red}{\mathbf{Red}}
\newcommand{\Proc}{\mathbf{Proc}}
\newcommand{\Infer}{\mathbf{Infer}}
\newcommand{\HNN}{\Hil_{NN}}
\newcommand{\HNNpi}{\Hil_{NN\pi}}
\newcommand{\StateBundle}{\texttt{N1NuclearStateBundle}}
\newcommand{\PionSector}{\texttt{NNPiSectorSpec}}
\newcommand{\ThreeBodyRecoil}{\texttt{ThreeBodyNuclearRecoilMap}}
\newcommand{\TransitionRecoil}{\texttt{NuclearTransitionRecoilMap}}
\newcommand{\PionOperator}{\texttt{PionActiveOperator}}
\newcommand{\TransitionOperator}{\texttt{NuclearTransitionOperator}}
\newcommand{\PionSubtraction}{\texttt{PionSubtractionMap}}
\newcommand{\TwoBodyCurrent}{\texttt{NuclearTwoBodyCurrent}}
\newcommand{\ContinuityLedger}{\texttt{NuclearContinuityLedger}}
\newcommand{\CoherentKernel}{\texttt{CoherentPropagationKernel}}
\newcommand{\CoherentBundle}{\texttt{CoherentAmplitudeBundle}}
\newcommand{\OverlapMap}{\texttt{PartonNuclearOverlapMap}}
\newcommand{\ProvComplex}{\texttt{Provenance2Complex}}
\newcommand{\AssumptionBundle}{\texttt{AssumptionBundle}}
\newcommand{\PredictionPlan}{\texttt{PredictionPlan}}

\newcolumntype{Y}{>{\raggedright\arraybackslash}X}
\newcolumntype{L}[1]{>{\raggedright\arraybackslash}p{#1}}

\title{\vspace{-1.0cm}
{\color{navy}\bfseries Volume XIII: Spin-Resolved \(NN\pi\) Dynamics, Pion Matching,\\[2mm]
Two-Body Currents, and Coherent Nuclear Amplitudes}\\[4mm]
\large Three-body light-front structure, pion-active and transition operators,\\
internal--exchange subtraction, finite-basis continuity, and a helicity-resolved small-\(x\) pilot}
\author{DeuteronWigner project}
\date{31 July 2026}

\begin{document}
\maketitle

\begin{center}
\begin{minipage}{0.94\textwidth}
\colorbox{softblue}{\parbox{0.96\textwidth}{
\textbf{Status and purpose.}
C15/N0 established the first isolated microscopic spin-1 nuclear validation
root: a normalized \(NN\) light-front state with S/D amplitudes, correlated
microscopic proton and neutron members, complete deuteron--parton helicity
parents, order-resolved partonic Wilson structure, common GTMD reductions,
consistent current benchmarks, and tagged-inclusive closure.  C16/N1 is the
active implementation package extending that state to \(NN\oplus NN\pi\).
The present volume is the normative formalism for that extension.  It treats
\(NN\pi\) as a spin-, isospin-, charge-, momentum-, and orbital-resolved
three-body amplitude; derives pion-active, nucleon-active, and transition
operator blocks from one sector theory; matches internal-nucleon and
inter-nucleon pion regions with an explicit overlap subtraction; constructs
currents consistently with the nuclear Hamiltonian; and introduces a
helicity-resolved coherent small-\(x\) amplitude without identifying it with
the partonic Wilson staple.  It does not claim a physical pion TMD, a complete
nuclear EFT, physical shadowing, a complete exchange-current basis, QCD
matching, evolution, process factorization, or inference readiness.
}}
\end{minipage}
\end{center}

\begin{abstract}
The first microscopic deuteron parent at the nucleonic level is not sufficient
to describe pion-active partonic structure, transition currents, or coherent
small-\(x\) dynamics.  These mechanisms must be represented by amplitudes and
operators in a common nuclear Hilbert space rather than by separately
normalized corrections.  We formulate a regulated two-sector light-front
model
\(
\Hil_{\mathrm{N1}}=\Hil_{NN}\oplus\Hil_{NN\pi}
\)
with one normalized spin-1 state, a Hermitian sector-changing Hamiltonian,
exact three-body recoil maps, charge-complete isospin coupling, parity and
orbital selection, diagonal nucleon-active and pion-active partonic operators,
and coherent \(NN\leftrightarrow NN\pi\) transition blocks.  A resolution map
separates pion-like configurations internal to the microscopic nucleon from
inter-nucleon exchange pions, producing a matched three-term combination and
an executable no-double-counting two-cell.  Electromagnetic, axial, and
energy--momentum currents are transformed with the same Hamiltonian and obey a
finite-resolution continuity ledger only when pion-in-flight, transition,
contact, induced, counterterm, and regulator pieces are included together.
The coherent small-\(x\) pilot is constructed from helicity-resolved
diffractive amplitudes and propagation phases; scalar, vector, and tensor
responses are projected only after amplitudes are combined.  Completely
positive reductions are derived by partial trace only after unresolved
coherence has been removed.  We give the full operator decomposition, tensor
network realization, provenance complex, benchmark suite, software
interfaces, and acceptance conditions for the N1 package and identify the
boundary to explicit \(\Delta\Delta\), compact six-quark, hidden-color, and
higher coherent sectors.
\end{abstract}

\tableofcontents

\section{Role, scope, and inherited contracts}
\label{sec:scope}

\subsection{The N0 result and the N1 transition}

C15/N0 supplies the normalized nucleonic state
\begin{equation}
 |D,\Lambda\rangle_{N0}=|\Psi_{NN}^{\Lambda}\rangle,
 \qquad
 \inner{\Psi_{NN}^{\Lambda}}{\Psi_{NN}^{\Lambda}}=1,
 \label{eq:n0-state}
\end{equation}
with one spectator-preserving recoil authority, the complete off-forward
nucleon spectral amplitude, spin-1 projectors, order-resolved microscopic
nucleon GTMD parents, current and angular-condition tests, and an explicit
nuclear tensor-network representation.  The N1 change of model space is
\begin{equation}
 \Hil_{N1}=\HNN\oplus\HNNpi,
 \label{eq:n1-space}
\end{equation}
with one normalized state
\begin{equation}
 |D,\Lambda\rangle_{N1}
 =|\Psi_{NN}^{\Lambda}\rangle
 +|\Psi_{NN\pi}^{\Lambda}\rangle .
 \label{eq:n1-state}
\end{equation}
The new sector is not an additive pion PDF placed on top of the N0 output.  It
changes the state, Hamiltonian, current, partonic operator, normalization,
momentum ledger, and provenance simultaneously.

\begin{decision}[N1 model class]
The first non-nucleonic nuclear extension is a regulated
\(NN\oplus NN\pi\) amplitude theory.  Pion-active, transition, exchange-current,
and coherent terms are different operator blocks of this theory.  They may not
be introduced as unrelated scalar corrections to named TMDs.
\end{decision}

\subsection{Inherited immutable identities}

Every N1 object inherits and preserves:
\begin{itemize}
\item the light-front convention
\(
 v^{\pm}=(v^0\pm v^3)/\sqrt2
\)
and the symmetric zero-skewness external fibers;
\item the distinct coordinates \(\bD\leftrightarrow\DT\) and
\(\bM\leftrightarrow\kT\);
\item correlated microscopic proton and neutron member identity;
\item Wilson order, future/past orientation, ordered gluon-link pair, and
independent \(f\)- and \(d\)-type color channels;
\item the spin-1 projector convention and the adapter
\(
 f_{1LL}=-(2/3)\delta_Tf_1
\);
\item separate amplitude, density, matching, reduction, and process map
classes;
\item the accepted production root and all phenomenological artifacts as
read-only regression oracles.
\end{itemize}

\begin{requirement}[No implicit promotion]
An N0 or N1 object may enter a physical TMD, process, evolution, or inference
calculation only through a separately validated matching or factorization map.
The fact that a regulated matrix element exists does not assign it a physical
QCD scheme.
\end{requirement}

\subsection{Primary N1 domain}

The primary N1 calculation is restricted to:
\begin{itemize}
\item zero skewness and finite transverse transfer;
\item the \(NN\) and \(NN\pi\) nuclear sectors;
\item one declared finite-resolution tower;
\item one Hermitian \(NN\leftrightarrow NN\pi\) transition interaction;
\item diagonal nucleon-active and pion-active partonic operators;
\item sector-transition partonic and current operators;
\item one analytic helicity-resolved coherent double-scattering pilot;
\item derived CP maps only after coherent amplitudes are combined;
\item exact, Krylov, full-bond, and reduced-bond numerical routes where the
state dimension permits them.
\end{itemize}

\subsection{Explicit nonclaims}

The strongest permissible N1 status is a finite-resolution validation status.
The present formalism does not establish:
\begin{itemize}
\item a physical deuteron prediction;
\item a physical pion GTMD, GPD, PDF, or TMD;
\item a complete chiral or nuclear EFT;
\item physical polarized or tensor nuclear shadowing;
\item a complete exchange-current basis;
\item nuclear Glauber or tagged final-state rescattering;
\item an explicit \(\Delta\Delta\), six-quark, hidden-color, or compact state;
\item LF-to-QCD ultraviolet matching, continuum soft subtraction, or a
physical TMD scheme;
\item Collins--Soper evolution, process hard factors, \(W+Y\), or inference.
\end{itemize}

\section{Resolution hierarchy and the two-sector nuclear theory}
\label{sec:resolution}

\subsection{Nuclear resolution identity}

A nuclear resolution is a typed tuple
\begin{equation}
 r_N=
 \left(
 \Lambda_N,N_{\max}^{N},K_N,
 \cF_N,
 \cB_N,
 \nu_J,
 \nu_I,
 \mathcal R_{\mathrm{end}},
 \mathcal R_{0},
 n_J,n_O,n_{\mathrm{coh}}
 \right),
 \label{eq:nuclear-resolution}
\end{equation}
where \(\cF_N\) lists the retained nuclear Fock sectors, \(\cB_N\) the basis,
\(\nu_J\) the current order, \(\nu_I\) the interaction order,
\(\mathcal R_{\mathrm{end}}\) the endpoint prescription,
\(\mathcal R_0\) the zero-mode policy, \(n_O\) the partonic-operator order,
and \(n_{\mathrm{coh}}\) the coherent-scattering order.  These labels are part
of every state and operator identity.

\subsection{Block Hamiltonian}

At level \(r_N\), the mass operator is
\begin{equation}
 \cM_{N1,r_N}^{2}
 =
 \begin{pmatrix}
 \cM_{NN,r_N}^{2} & V_{NN\leftarrow NN\pi,r_N}\\
 V_{NN\pi\leftarrow NN,r_N} & \cM_{NN\pi,r_N}^{2}
 \end{pmatrix},
 \qquad
 V_{NN\leftarrow NN\pi}=V_{NN\pi\leftarrow NN}^{\dagger}.
 \label{eq:n1-hamiltonian}
\end{equation}
The diagonal \(NN\) block is the N0 Hamiltonian.  The \(NN\pi\) block contains
three-body kinetic energy, pion and nucleon masses, retained interactions,
induced terms, and sector-specific counterterms.  The off-diagonal block
creates or absorbs the pion while preserving total charge, baryon number,
\(I\), \(J^P\), \(J^z\), momentum, regulator, and phase conventions.

\begin{requirement}[Generated adjoint]
The absorption block is generated from the emission block as the Hilbert-space
adjoint.  Independent parameterizations of the two directions are forbidden.
\end{requirement}

\subsection{Renormalization trajectory}

The N1 renormalization datum is
\begin{equation}
 \mathfrak R_{N1,r}
 =
 \left(
 \mathcal R_r,
 \theta_{NN,r}^{\bare},
 \theta_{NN\pi,r}^{\bare},
 \delta\theta_{NN,r},
 \delta\theta_{NN\pi,r},
 \delta g_{NN\pi,r},
 \{\mathcal C_i\},
 \mathcal S_r,
 \Delta_r
 \right).
 \label{eq:n1-renormalization}
\end{equation}
The same physical calibration conditions are imposed at every resolution,
while bare parameters and counterterms flow.  At minimum the calibration set
contains the deuteron pole, charge, one transition condition, and one current
condition.  Tensor, coherent, pion-active, and second-current-component
observables remain holdouts.

\subsection{State normalization and plus-momentum ledger}

Define
\begin{equation}
 Z_{NN}=\inner{\Psi_{NN}}{\Psi_{NN}},
 \qquad
 Z_{NN\pi}=\inner{\Psi_{NN\pi}}{\Psi_{NN\pi}},
 \qquad
 Z_{NN}+Z_{NN\pi}=1.
 \label{eq:sector-normalization}
\end{equation}
The plus-momentum ledger is
\begin{equation}
 1=
 \sum_{\alpha\in\{NN,NN\pi\}}
 \int[\dd\Gamma_{\alpha}]
 \left(\sum_{i\in\alpha}y_i\right)
 \abs{\Psi_{\alpha}}^2.
 \label{eq:n1-plus-ledger}
\end{equation}
Individual sector probabilities are resolution dependent.  Total
normalization, charge, baryon number, isospin, and matched momentum moments are
the physical closure conditions.

\section{Spin, isospin, parity, and statistics of the \texorpdfstring{\(NN\pi\)}{NNpi} sector}
\label{sec:sector}

\subsection{Exact quantum numbers}

At the charge-symmetric reference point, the sector has
\begin{equation}
 B=2,
 \qquad Q=+1,
 \qquad I=0,
 \qquad J^P=1^+.
 \label{eq:nnpi-quantum-numbers}
\end{equation}
Because \(I_\pi=1\), the \(NN\) subsystem in the \(NN\pi\) sector must carry
\(I_{NN}=1\) in order to couple to total \(I=0\).  Charge-symmetry-breaking
admixtures are represented by separate assumption branches.

\subsection{Charge-complete isospin coupling}

The complete total-charge-one basis contains
\begin{equation}
 |pn\pi^0\rangle,
 \qquad
 |pp\pi^-\rangle,
 \qquad
 |nn\pi^+\rangle.
 \label{eq:charge-basis}
\end{equation}
Using a normalized \(I_{NN}=1\) basis, the total-isospin-zero state is
\begin{equation}
 |I=0,I_z=0\rangle
 =\frac{1}{\sqrt3}
 \left(
 |1,+1\rangle_{NN}|1,-1\rangle_{\pi}
 -|1,0\rangle_{NN}|1,0\rangle_{\pi}
 +|1,-1\rangle_{NN}|1,+1\rangle_{\pi}
 \right).
 \label{eq:isospin-zero-state}
\end{equation}
The phase convention is versioned and tested.  Retaining only \(pn\pi^0\)
violates isospin completeness.

\subsection{Parity and angular momentum}

Let \(\ell_{NN}\) denote the orbital angular momentum inside the nucleon pair
and \(\ell_{\pi}\) the Jacobi orbital momentum of the pion relative to the
pair.  Since the pion has intrinsic parity \(-1\),
\begin{equation}
 P_{NN\pi}
 =(-1)^{\ell_{NN}+\ell_{\pi}}(-1).
 \label{eq:nnpi-parity}
\end{equation}
The deuteron channel therefore requires
\begin{equation}
 \ell_{NN}+\ell_{\pi}\quad\text{odd}.
 \label{eq:odd-orbital-parity}
\end{equation}
The minimal benchmark basis must contain the lowest odd total-orbital
combinations capable of coupling the pair spin and pion orbital motion to
\(J=1\).  Higher partial waves form a controlled truncation sequence.

\subsection{Nucleon exchange statistics}

The two nucleons are components of one isospin doublet and obey fermionic
exchange statistics.  For the nucleon pair,
\begin{equation}
 (-1)^{\ell_{NN}+S_{NN}+I_{NN}}=-1.
 \label{eq:pauli-nnpi}
\end{equation}
With \(I_{NN}=1\), the retained \((\ell_{NN},S_{NN})\) channels must satisfy
this rule.  The \(pp\) and \(nn\) charge channels are not exempt from the same
antisymmetry constraint.

\begin{requirement}[Amplitude, not probability]
The \(NN\pi\) sector is represented by spin-, isospin-, charge-, orbital-, and
momentum-resolved amplitudes.  A scalar pion probability cannot define a
tensor-polarized GTMD or transition current.
\end{requirement}

\section{Three-body light-front kinematics and typed recoil}
\label{sec:kinematics}

\subsection{External fibers and intrinsic variables}

Use the symmetric zero-skewness deuteron frame
\begin{equation}
 P_i=P-\frac{\Delta}{2},
 \qquad
 P_f=P+\frac{\Delta}{2},
 \qquad
 \Delta^+=0,
 \qquad
 P_{iT}=-\frac{\DT}{2},
 \quad
 P_{fT}=+\frac{\DT}{2}.
 \label{eq:n1-symmetric-frame}
\end{equation}
For the three-body sector introduce positive fractions
\begin{equation}
 y_1>0,
 \qquad y_2>0,
 \qquad y_\pi>0,
 \qquad y_1+y_2+y_\pi=1,
 \label{eq:threebody-fractions}
\end{equation}
and average intrinsic transverse momenta
\begin{equation}
 \bm\kappa_{1T}+\bm\kappa_{2T}+\bm\kappa_{\pi T}=0.
 \label{eq:threebody-closure}
\end{equation}
A convenient invariant measure is
\begin{equation}
 [\dd\Gamma_3]
 =
 \frac{\dd y_1\dd y_2\dd y_\pi\,
 \delta(1-y_1-y_2-y_\pi)}
 {2(2\pi)^6 y_1y_2y_\pi}
 \prod_{i=1,2,\pi}\dd^2\bm\kappa_{iT}
 \delta^{(2)}\!\left(\sum_i\bm\kappa_{iT}\right),
 \label{eq:threebody-measure}
\end{equation}
up to a declared unitary normalization adapter.  The wave function, current,
operator, and convolution Jacobians transform together under any alternative
normalization.

\subsection{Jacobi coordinates}

One useful Jacobi choice is
\begin{align}
 \alpha&=\frac{y_1}{y_1+y_2},
 &
 \bm r_T&=
 \frac{y_2\bm\kappa_{1T}-y_1\bm\kappa_{2T}}{y_1+y_2},
 \label{eq:nn-jacobi}\\
 \beta&=y_\pi,
 &
 \bm s_T&=\bm\kappa_{\pi T}.
 \label{eq:pion-jacobi}
\end{align}
The inverse map and Jacobian are part of \PionSector.  Another Jacobi tree is
permitted only through an explicit recoupling adapter.

\subsection{Diagonal active-constituent recoil}

Let constituent \(j\in\{1,2,\pi\}\) be active.  Define
\begin{align}
 \bm\kappa_{jT}^{\mathrm{in}}
 &=\bm k_{jT}-\frac{1-y_j}{2}\DT,
 &
 \bm\kappa_{jT}^{\mathrm{out}}
 &=\bm k_{jT}+\frac{1-y_j}{2}\DT,
 \label{eq:active-threebody-recoil}\\
 \bm\kappa_{iT}^{\mathrm{in}}
 &=\bm k_{iT}+\frac{y_i}{2}\DT,
 &
 \bm\kappa_{iT}^{\mathrm{out}}
 &=\bm k_{iT}-\frac{y_i}{2}\DT,
 \qquad i\neq j.
 \label{eq:spectator-threebody-recoil}
\end{align}

\begin{proposition}[Three-body spectator theorem]
With \cref{eq:active-threebody-recoil,eq:spectator-threebody-recoil}, the
intrinsic momenta close in both fibers, the active physical constituent
receives the full transfer \(\Delta\), every physical spectator momentum is
unchanged, and the transformation has unit Jacobian.
\end{proposition}

\begin{proof}
Intrinsic closure follows from
\(
-(1-y_j)+\sum_{i\neq j}y_i=0
\).
Writing the physical transverse momentum as
\(
 p_{iT}=y_iP_T+\kappa_{iT}
\), one obtains
\(
 p_{jT}^{\mathrm{out}}-p_{jT}^{\mathrm{in}}=\DT
\)
and
\(
 p_{iT}^{\mathrm{out}}=p_{iT}^{\mathrm{in}}
\)
for every spectator.  The map is a translation and therefore has unit
Jacobian.
\end{proof}

\subsection{Transition recoil for \texorpdfstring{\(NN\leftrightarrow NN\pi\)}{NN to NNpi}}

A particle-number-changing operator requires a distinct map
\begin{equation}
 \TransitionRecoil:
 \cM_{NN}^{\mathrm{in}}
 \times
 \cO_{NN\to NN\pi}
 \longrightarrow
 \cM_{NN\pi}^{\mathrm{out}},
 \label{eq:transition-recoil-map}
\end{equation}
with inverse for pion absorption.  It stores:
\begin{itemize}
\item the emitting nucleon;
\item pion fraction and transverse momentum;
\item recoil sharing among the residual pair;
\item external transfer routing;
\item endpoint and zero-mode prescriptions;
\item the Jacobian and normalization convention;
\item source and target sector identities.
\end{itemize}
A diagonal three-body recoil map cannot be reused by changing only an array
shape.

\subsection{Support and endpoint policy}

The regulated support is
\begin{equation}
 y_i\geq y_{\min}(r_N)>0,
 \qquad
 \sum_i y_i=1,
 \label{eq:endpoint-policy}
\end{equation}
with a declared limit as the longitudinal resolution is enlarged.  The
endpoint regulator is part of every pion-active and transition-operator
identity.  A model result is not endpoint converged merely because a
quadrature at one \(y_{\min}\) is stable.

\section{The spin-resolved \texorpdfstring{\(NN\pi\)}{NNpi} light-front amplitude}
\label{sec:amplitude}

\subsection{Amplitude expansion}

The three-body state is
\begin{equation}
 \begin{aligned}
 |\Psi_{NN\pi}^{\Lambda}\rangle
 ={}&
 \sum_{c,\lambda_1,\lambda_2}
 \sum_{\alpha_{I},\alpha_{L},\alpha_{S}}
 \int[\dd\Gamma_3]
 \,\Psi_{c;\lambda_1\lambda_2;
 \alpha_I\alpha_L\alpha_S}^{\Lambda}
 (y_i,\bm\kappa_{iT})
 \\
 &\times
 |N_1N_2\pi;c,\lambda_1,\lambda_2,
 \alpha_I,\alpha_L,\alpha_S;
 y_i,\bm\kappa_{iT}\rangle,
 \end{aligned}
 \label{eq:nnpi-amplitude}
\end{equation}
where \(c\) labels the charge channel and the \(\alpha\)'s retain isospin,
orbital, spin, and multiplicity information.

\subsection{Coupling grammar}

A symmetry-adapted fusion tree may be written schematically as
\begin{equation}
 \left[
 \left(N_1\otimes N_2\right)_{S_{NN},I_{NN},\ell_{NN}}
 \otimes
 \pi_{I_\pi=1,\ell_\pi}
 \right]_{J=1,I=0,P=+}.
 \label{eq:nnpi-fusion-tree}
\end{equation}
Every Clebsch--Gordan coefficient, phase convention, outer multiplicity, and
fermionic exchange sign is stored rather than reconstructed from a final
channel name.

\subsection{State generation and comparison routes}

The central N1 state may be generated by:
\begin{enumerate}
\item an explicit diagonalization of \cref{eq:n1-hamiltonian};
\item a matrix-free Krylov solver;
\item exact full-bond tensorization;
\item a variational nuclear TTN;
\item an analytic weak-coupling oracle used only for validation.
\end{enumerate}
Exact and full-bond states must agree.  Reduced-bond states are compared on
pion probability, transition interference, tensor observables, currents, and
coherent amplitudes, not only on the eigenvalue.

\subsection{Controlled relation to light-front Hamiltonian EFT}

Light-front Hamiltonian EFT calculations with explicit pion Fock sectors
provide an important feasibility route for nonperturbative pion distributions
and sea-flavor asymmetry \cite{CaoEtAl2026}.  They do not, by themselves,
supply the complete spin-resolved transverse \(NN\pi\) amplitude required here
unless their state, recoil, current, and operator information is exported with
the present identities.  N1 may use such a calculation as a future matched
input or comparison, not as an unnamed replacement for the three-body state.

\section{Matched N1 partonic operator basis}
\label{sec:operators}

\subsection{Sector matrix}

For parton species \(a\), Dirac or Lorentz projection \(J\), and Wilson identity
\(\gamma\), define
\begin{equation}
 \widehat\cO_{a}^{[J,\gamma],N1}
 =
 \begin{pmatrix}
 \widehat\cO_{a,NN\to NN}
 &
 \widehat\cO_{a,NN\leftarrow NN\pi}
 \\
 \widehat\cO_{a,NN\pi\leftarrow NN}
 &
 \widehat\cO_{a,NN\pi\to NN\pi}
 \end{pmatrix}.
 \label{eq:n1-operator-matrix}
\end{equation}
Hermiticity requires the transition blocks to be adjoints after transfer,
path, and helicity reversal.

\subsection{Diagonal \texorpdfstring{\(NN\pi\)}{NNpi} blocks}

The diagonal block decomposes as
\begin{equation}
 \widehat\cO_{a,NN\pi\to NN\pi}
 =
 \widehat\cO_{a,N_1}^{(1)}
 +\widehat\cO_{a,N_2}^{(1)}
 +\widehat\cO_{a,\pi}^{(1)}
 +\widehat\cO_{a,NN\pi}^{(2)}
 +\widehat\cO_{a,NN\pi}^{\ind}.
 \label{eq:nnpi-diagonal-operator}
\end{equation}
The first two terms describe an active nucleon with the other nucleon and pion
as spectators.  The third is the pion-active operator.  The remaining terms
are matched two-body and induced operators.

\subsection{Transition operators}

The sector-changing partonic operator is
\begin{equation}
 \widehat\cO_{a,NN\pi\leftarrow NN}
 =
 \widehat\cO_{a,\mathrm{emit}}
 +\widehat\cO_{a,\mathrm{contact}}
 +\widehat\cO_{a,\mathrm{ind}}
 +\delta\widehat\cO_{a,\ct},
 \label{eq:transition-operator}
\end{equation}
with the reverse block generated as its adjoint.  A transition operator can
carry partonic transfer through the active nucleon, the emitted pion, or a
contact block.  Each routing is a distinct typed contribution.

\subsection{Complete operator identity}

Every N1 operator stores at least
\begin{equation}
 \begin{aligned}
 \mathfrak o_{a/D}^{N1}=
 (&a,J,P_i,P_f,\xi,\Delta,
 \text{source sector},\text{target sector},
 \text{active constituent},
 \\
 &\gamma,R_c,\text{ordered gluon links},f/d,
 n_W,\mathcal R_{\UV},\mathcal R_{\mathrm{rap}},
 \mathcal S_{\mathrm{soft}},
 \\
 &\mu,\zeta,\mathfrak s,
 r_N,\mathfrak m_N,\mathfrak m_D,
 \text{recoil ID},\text{current family}).
 \end{aligned}
 \label{eq:n1-operator-id}
\end{equation}
An operator with missing sector, active-constituent, recoil, path, or member
identity cannot enter the common parent.

\subsection{Pion-active operator status}

The initial pion partonic operator may be a regulated analytic validation
oracle.  It retains pion charge, isospin, direct positive-\(x\) quark and
antiquark identity, gluon status, number and momentum ledgers, source hash,
and matching limitations.  It may not be labeled a physical pion GTMD or TMD,
and its normalization may not be fitted independently to a deuteron
observable.

\section{The full N1 deuteron GTMD parent}
\label{sec:parent}

\subsection{Amplitude-level definition}

The common parent is
\begin{equation}
 \begin{aligned}
 W_{a/D}^{[J,\gamma],N1}
 ={}&
 \langle\Psi_{NN}|\widehat\cO_{NN\to NN}|\Psi_{NN}\rangle
 \\
 &+
 \langle\Psi_{NN\pi}|\widehat\cO_{NN\pi\to NN\pi}|\Psi_{NN\pi}\rangle
 \\
 &+
 \langle\Psi_{NN\pi}|\widehat\cO_{NN\pi\leftarrow NN}|\Psi_{NN}\rangle
 \\
 &+
 \langle\Psi_{NN}|\widehat\cO_{NN\leftarrow NN\pi}|\Psi_{NN\pi}\rangle
 \\
 &+\delta W_{a/D}^{\mathrm{coh,pilot}},
 \end{aligned}
 \label{eq:n1-parent}
\end{equation}
where the final term is present only in an explicitly selected coherent-pilot
branch.  The first four terms form the exact parent of the two-sector state.

\begin{requirement}[Project after composition]
The target channels \(U,L,T,LL,LT,TT\), named TMDs, and collinear moments are
projected only after all selected sector-diagonal and transition amplitudes
have been combined.  Sector-specific scalar functions may not be normalized
independently and added after projection.
\end{requirement}

\subsection{Hermiticity and discrete symmetries}

At zero skewness,
\begin{equation}
 W_{\Lambda'\Lambda}(x,\kT,\DT;\gamma)
 =
 W_{\Lambda\Lambda'}^{*}
 (x,\kT,-\DT;\gamma^{-1}),
 \label{eq:n1-hermiticity}
\end{equation}
with the appropriate target, parton, charge, and sector phases.  Parity acts
on the complete orbital and pion-parity structure.  Time reversal maps the
partonic Wilson paths but does not, by itself, identify nuclear coherent
propagation with the partonic staple.

\subsection{Spin-1 and parton-helicity matrices}

For quarks and antiquarks the parent is stored as a \(6\times6\) joint matrix
in deuteron and parton helicities.  For gluons it is stored both as a target
matrix with transverse field indices and as an equivalent target--gluon
helicity matrix.  The nine target projectors reconstruct an arbitrary
\(3\times3\) matrix.  The convention-independent tensor combination is
\begin{equation}
 \delta_TF
 =F_{\Lambda=0}
 -\frac12\left(F_{\Lambda=+1}+F_{\Lambda=-1}\right),
 \label{eq:n1-deltaT}
\end{equation}
with the project adapter
\begin{equation}
 f_{1LL}=-\frac23\,\delta_Tf_1.
 \label{eq:n1-f1ll}
\end{equation}

\subsection{Partonic Wilson identity through nuclear composition}

The N1 parent preserves the nucleon-side Wilson order and path data.  In
particular, gluon ordered-link pairs and independent \(f\)- and \(d\)-type
channels survive composition.  The following remain distinct:
\begin{equation}
 \texttt{PARTONIC\_WILSON\_STAPLE}
 \neq
 \texttt{NUCLEAR\_COHERENT\_PROPAGATION}
 \neq
 \texttt{TAGGED\_FSI}.
 \label{eq:three-rescattering-types}
\end{equation}
No nuclear coherent phase may be inserted by modifying the partonic path
identity.

\subsection{Common-member covariance}

One N1 member contains the correlated proton and neutron microscopic states,
nuclear state, pion sector, transition interaction, current family, partonic
operator, Wilson identity, and numerical realization.  Observable
covariances are evaluated member by member.  A proton from one microscopic
plan and a neutron or pion sector from another cannot form one nuclear member.

\section{Pion-active and transition reductions}
\label{sec:pion-reductions}

\subsection{Pion-active convolution as a reduction}

When the probe acts on the pion and the two nucleons are unresolved
spectators, the diagonal contribution reduces to
\begin{equation}
 \begin{aligned}
 W_{a/D}^{(\pi)}(x_D,\kT,\DT)
 ={}&
 \sum_{c_\pi}
 \int_{x_D}^{1}\frac{\dd y_\pi}{y_\pi}
 \int\frac{\dd^2\qT}{(2\pi)^2}
 \\
 &\times
 \rho_{\pi/D}^{c_\pi}
 (y_\pi,\qT,\DT)
 W_{a/\pi^{c_\pi}}
 \left(
 \frac{x_D}{y_\pi},
 \kT-\frac{x_D}{y_\pi}\qT,
 \DT
 \right)
 \cJ_{\pi}.
 \end{aligned}
 \label{eq:pion-convolution}
\end{equation}
This is a Sullivan-type reduction of the complete \(NN\pi\) matrix element,
not an independent postulate \cite{Sullivan1972,Miller2014}.

\subsection{Nucleon-active terms in the \texorpdfstring{\(NN\pi\)}{NNpi} sector}

For an active nucleon with a spectator nucleon and pion,
\begin{equation}
 \begin{aligned}
 W_{a/D}^{(N|N\pi)}
 ={}&
 \sum_{N,c_\pi}
 \int[\dd\Gamma_3]
 \sum_{\lambda',\lambda}
 \rho_{\Lambda'\Lambda;\lambda'\lambda}^{N/(NN\pi)}
 \\
 &\times
 W_{a/N,\lambda'\lambda}
 \left(
 \frac{x_D}{y_N},
 \kT-\frac{x_D}{y_N}\bm\kappa_{NT},
 \DT
 \right)
 \cJ_{N|N\pi}.
 \end{aligned}
 \label{eq:nucleon-active-nnpi}
\end{equation}
The microscopic nucleon parent and the explicit nuclear pion sector must obey
the pion-overlap subtraction of \cref{sec:subtraction}.

\subsection{Transition interference}

The off-diagonal contribution is
\begin{equation}
 W_{a/D}^{\mathrm{trans}}
 =
 2\operatorname{Re}
 \left[
 \langle\Psi_{NN\pi}|
 \widehat\cO_{a,NN\pi\leftarrow NN}
 |\Psi_{NN}\rangle
 \right],
 \label{eq:transition-interference}
\end{equation}
for Hermitian T-even operators, with the corresponding matrix-valued relation
for general helicity and Wilson structures.  Transition interference can
contribute to tensor observables even when the diagonal pion probability is
small.

\subsection{TMD, GPD, PDF, and current closure}

The N1 parent must satisfy the regulated commuting reductions
\begin{align}
 \mathsf T[W]&=W(x,\kT,\DT=0),
 \label{eq:n1-tmd-reduction}\\
 \mathsf G[W]&=\int\dd^2\kT\,W(x,\kT,\DT),
 \label{eq:n1-gpd-reduction}\\
 \mathsf P[W]&=\left.\int\dd^2\kT\,W\right|_{\DT=0},
 \label{eq:n1-pdf-reduction}\\
 \mathsf Q[W]&=\int\dd x\,\dd^2\kT\,W.
 \label{eq:n1-current-reduction}
\end{align}
Direct and sequential routes must agree for the full sector sum and for each
provenance-resolved block.  The physical QCD interpretation remains gated by
operator matching and soft/rapidity renormalization.

\subsection{Tensor distribution and \texorpdfstring{\(b_1\)}{b1}}

The full N1 tensor PDF is
\begin{equation}
 \delta_T f_{a/D}^{N1}
 =
 \delta_T f_{a/D}^{NN}
 +\delta_T f_{a/D}^{NN\pi,\mathrm{diag}}
 +\delta_T f_{a/D}^{NN\leftrightarrow NN\pi}
 +\delta_T f_{a/D}^{\mathrm{coh,pilot}},
 \label{eq:n1-tensor-pdf}
\end{equation}
with only selected branches present.  At leading order,
\begin{equation}
 b_1(x,Q^2)
 =\frac12\sum_q e_q^2
 \left[
 \delta_Tq_D(x,Q^2)+\delta_T\bar q_D(x,Q^2)
 \right]
 +\cO(\alpha_s,M_D^2/Q^2).
 \label{eq:n1-b1}
\end{equation}
The pion-active and transition components are predictions of the common N1
state and operator.  They do not receive independent \(b_1\) normalizations.

\section{Internal-nucleon versus nuclear-exchange pion matching}
\label{sec:subtraction}

\subsection{The overlap problem}

The microscopic H7 nucleon already contains pion-like light-sea and chiral
amplitudes.  The N1 state additionally contains an explicit pion exchanged or
propagating between nucleons.  Without a separation prescription, the same
low-resolution pion region can be counted twice.

Define typed descriptions:
\begin{equation}
 \mathfrak P_{\mathrm{int}},
 \qquad
 \mathfrak P_{\mathrm{exch}},
 \qquad
 \mathfrak P_{\mathrm{ind}},
 \label{eq:pion-description-types}
\end{equation}
for internal-nucleon, explicit exchange-sector, and induced/contact
representations.

\subsection{Resolution separator}

Introduce a smooth or sharp projector
\begin{equation}
 \cP_{\pi}^{\mathrm{int}}(\omega),
 \qquad
 \cP_{\pi}^{\mathrm{exch}}(\omega),
 \qquad
 \cP_{\pi}^{\mathrm{int}}+\cP_{\pi}^{\mathrm{exch}}
 =1+\delta_{\omega},
 \label{eq:pion-separator}
\end{equation}
where \(\omega\) denotes the declared resolution variable and
\(\delta_{\omega}\) measures truncation or nonorthogonality.  The separator is
not a probability observable; it is a matching convention.

\subsection{Matched pion contribution}

The matched contribution is
\begin{equation}
 W_{\pi}^{\mathrm{matched}}
 =
 W_{\pi}^{\mathrm{internal}}
 +W_{\pi}^{\mathrm{exchange}}
 -W_{\pi}^{\mathrm{overlap}}.
 \label{eq:pion-matched-sum}
\end{equation}
Changing the separation convention may redistribute the three terms, but the
matched sum must remain stable within the declared truncation uncertainty.

\begin{proposition}[Pion representation covariance]
Suppose two pion-resolution conventions are related by an invertible
model-space transformation and all Hamiltonian and partonic operators are
transformed consistently.  Then the complete matched matrix element is
invariant to the retained order, while the internal, exchange, and overlap
components are representation dependent.
\end{proposition}

\subsection{Explicit versus induced pion theory}

Eliminating \(NN\pi\) through a Feshbach map gives
\begin{equation}
 H_{\eff}(E)
 =PHP+PHQ(E-QHQ)^{-1}QHP,
 \label{eq:n1-feshbach-h}
\end{equation}
and transforms any operator as
\begin{equation}
 O_{\eff}(E',E)
 =P[1+\omega^{\dagger}(E')]O[1+\omega(E)]P.
 \label{eq:n1-feshbach-o}
\end{equation}
The provenance relation is
\begin{equation}
 \texttt{EXPLICIT\_NNPI}
 \Longleftrightarrow
 \texttt{INDUCED\_PION\_OPERATOR}
 +\texttt{TRANSFORMED\_OBSERVABLES}
 +\texttt{REMAINDER}.
 \label{eq:pion-provenance-equivalence}
\end{equation}
It is never an additive relation.

\subsection{No-double-counting two-cells}

The provenance complex contains two-cells for:
\begin{itemize}
\item internal plus exchange minus overlap;
\item explicit sector versus induced/contact replacement;
\item pion-active versus transition routing;
\item partonic soft overlap versus nuclear coherent overlap;
\item count-once equivalent cut descriptions.
\end{itemize}
Removing a required subtraction or selecting a representation together with
its replacement fails before numerical evaluation.

\section{Hamiltonian-consistent one- and two-body currents}
\label{sec:currents}

\subsection{Current family}

The N1 electromagnetic current is
\begin{equation}
 J_{N1}^{\mu}
 =
 J_N^{\mu}
 +J_{\pi}^{\mu}
 +J_{\mathrm{trans}}^{\mu}
 +J_{\mathrm{contact}}^{\mu}
 +J_{\mathrm{ind}}^{\mu}
 +\delta J_{\ct}^{\mu}.
 \label{eq:n1-current}
\end{equation}
The terms denote nucleon one-body, pion-in-flight, sector-transition,
contact/seagull, induced, and counterterm contributions.  Each is owned by the
same Hamiltonian identity as \cref{eq:n1-hamiltonian}.  Chiral EFT derivations
provide a systematic example of how one- and two-body current operators are
ordered and renormalized consistently with nuclear interactions
\cite{Baroni2015,Krebs2020,Schiavilla2018}.

The energy--momentum tensor and axial current have analogous decompositions,
with operator-specific contact and pion terms.  A current from an unrelated
Hamiltonian is inadmissible even if it reproduces the same static charge.

\subsection{Finite-resolution continuity identity}

The validation identity is
\begin{equation}
 q_{\mu}J_{N1}^{\mu}
 =
 [H_{N1},\rho_{N1}]
 +\delta_{\mathrm{cont}},
 \label{eq:n1-continuity}
\end{equation}
where the residual is decomposed as
\begin{equation}
 \delta_{\mathrm{cont}}
 =
 \delta_N
 +\delta_{\pi}
 +\delta_{\mathrm{trans}}
 +\delta_{\mathrm{contact}}
 +\delta_{\mathrm{ind}}
 +\delta_{\ct}
 +\delta_{\reg}
 +\delta_{\mathrm{basis}}.
 \label{eq:continuity-ledger}
\end{equation}
Omitting a required nonzero contribution must give a signed defect.  Passing
this benchmark permits only
\texttt{FINITE\_BASIS\_NUCLEAR\_CONTINUITY\_BENCHMARKED}.

\subsection{Spin-1 current and angular condition}

The plus-current helicity amplitudes
\begin{equation}
 I_{\Lambda'\Lambda}^{+}(Q^2)
 =\langle P',\Lambda'|J^{+}_{N1}|P,\Lambda\rangle
 \label{eq:n1-current-helicity}
\end{equation}
are reduced to charge, magnetic, and quadrupole form factors through a
versioned convention adapter.  The light-front angular condition is evaluated
with the complete current family.  It may not be repaired by assigning a
separate normalization to each helicity amplitude.

\subsection{Pion and exchange-current moments}

The direct current and integrated GTMD routes must agree:
\begin{equation}
 \int\dd x_D\,\dd^2\kT\,
 W_{a/D}^{[V],N1}
 \longleftrightarrow
 \langle D'|J_a^{+}|D\rangle,
 \label{eq:n1-current-gtmd-closure}
\end{equation}
with separate ledgers for nucleon one-body, pion-active, transition, contact,
and induced terms.  Exchange-current contributions to deuteron energy--momentum
form factors provide a useful independent benchmark of current conservation
and pion/seagull structure \cite{HeZahed2024}.

\subsection{Axial and PCAC extension}

When the N1 axial sector is enabled,
\begin{equation}
 A_{N1}^{\mu,a}
 =A_N^{\mu,a}+A_{\pi}^{\mu,a}
 +A_{\mathrm{trans}}^{\mu,a}
 +A_{\mathrm{contact}}^{\mu,a}
 +A_{\mathrm{ind}}^{\mu,a}
 +\delta A_{\ct}^{\mu,a},
 \label{eq:n1-axial-current}
\end{equation}
with a matched pseudoscalar operator.  The finite-basis PCAC residual is kept
separate from the electromagnetic continuity residual.  An explicit nuclear
pion sector does not automatically replace the induced pion-pole operator in
the microscopic nucleon; their overlap is controlled by the same pion
separator as \cref{sec:subtraction}.

\section{Helicity-resolved coherent small-\texorpdfstring{\(x\)}{x} amplitudes}
\label{sec:coherent}

\subsection{Amplitude-level origin}

The leading N1 coherent pilot is a two-step amplitude
\begin{equation}
 \delta W_{\Lambda'\Lambda}^{\mathrm{coh}}
 =
 \sum_X
 \sum_{\lambda_1,\lambda_2}
 \mathcal A_{\Lambda'\to X;\lambda_2}^{\mathrm{diff}\,*}
 \,\mathcal G_X
 \,\mathcal A_{\Lambda\to X;\lambda_1}^{\mathrm{diff}}
 \,\mathcal S_{\lambda_2\lambda_1}^{D},
 \label{eq:coherent-amplitude}
\end{equation}
where \(X\) is the diffractive intermediate state, \(\mathcal G_X\) carries
the propagation phase and longitudinal ordering, and \(\mathcal S^D\) is the
spin-resolved nuclear overlap.  This structure follows the leading-twist
connection between diffraction and nuclear shadowing, but the N1 pilot does
not claim the full phenomenology of that framework \cite{FrankfurtGuzeyStrikman2012}.

\subsection{Diffractive helicity identity}

Each elementary amplitude stores
\begin{equation}
 \mathfrak a_X=
 \left(
 a,J,x,Q^2,t,
 \lambda_N^{\mathrm{in}},\lambda_N^{\mathrm{out}},
 \lambda_X,
 \text{target constituent},
 \text{phase},
 \text{normalization},
 \text{source}
 \right).
 \label{eq:diff-amplitude-id}
\end{equation}
Unpolarized diffraction does not determine vector- or tensor-polarized
diffractive amplitudes.  The tensor coherent correction must be projected from
explicit helicity amplitudes, not copied from an unpolarized ratio.

\subsection{Propagation phase and ordering}

For longitudinal positions \(z_1,z_2\), a schematic propagator is
\begin{equation}
 \mathcal G_X(z_2,z_1)
 =\theta(z_2-z_1)
 \exp\left[i\Delta_X(z_2-z_1)\right]
 \mathcal A_X^{\mathrm{atten}},
 \label{eq:coherent-propagator}
\end{equation}
with a declared coherence length and attenuation status.  Reversing the
longitudinal ordering produces the corresponding conjugate amplitude.  Zero
elementary diffraction or zero propagation length must remove the correction.

\subsection{Scalar, vector, and tensor projections}

After amplitudes are combined, the target matrix is projected onto
\(K=0,1,2\) irreducible sectors.  A copied scalar shadowing factor would imply
\begin{equation}
 \delta W^{(2)}_{\mathrm{tensor}}
 \propto
 \delta W^{(0)}_{\mathrm{scalar}},
 \label{eq:copied-shadowing-ansatz}
\end{equation}
which the benchmark is required to falsify whenever the elementary helicity
amplitudes are nontrivial.

\subsection{Shadowing and antishadowing ledgers}

The coherent correction and any antishadowing compensation are distinct typed
mechanisms.  Antishadowing is constrained through the appropriate number,
momentum, and tensor moment ledgers; it is not an arbitrary enhancement
attached to every TMD.

\subsection{Partonic--nuclear overlap subtraction}

The partonic Wilson staple and nuclear coherent propagation may share a soft
or Glauber-sensitive region.  Introduce
\begin{equation}
 W^{\mathrm{matched}}
 =W^{\mathrm{partonic}}
 +W^{\mathrm{nuclear\ coh}}
 -W^{\mathrm{parton\text{-}nuclear\ overlap}}.
 \label{eq:parton-nuclear-subtraction}
\end{equation}
The overlap term carries path, scattering order, color, rapidity, nuclear
member, and resolution identity.  Missing and duplicate subtraction must
produce signed residuals.  N1 does not claim a complete continuum Glauber
factorization proof.

\section{Amplitude-level coherence and completely positive reductions}
\label{sec:cp}

\subsection{Amplitude embedding}

Let
\begin{equation}
 V:
 \Hil_{\mathrm{resolved}}
 \longrightarrow
 \Hil_D\otimes\Hil_{\mathrm{unresolved}}
 \label{eq:amplitude-embedding}
\end{equation}
be the full amplitude embedding after the \(NN\), \(NN\pi\), transition, and
coherent contributions have been combined.  A reduced map is then
\begin{equation}
 \cE(\rho)
 =\Tr_{\mathrm{unresolved}}
 \left[V\rho V^{\dagger}\right].
 \label{eq:derived-cp-map}
\end{equation}
This map is completely positive by construction.

\subsection{Where the reduction is valid}

The CP description is valid when the traced subsystem is not resolved by the
observable and all amplitudes leading to the same retained state have already
been summed.  It is useful for inclusive probability closure, detector
coarse-graining, and selected unresolved spectator channels.

\subsection{Premature trace failure}

If \(V=V_{NN}+V_{NN\pi}\), then
\begin{equation}
 V\rho V^{\dagger}
 =V_{NN}\rho V_{NN}^{\dagger}
 +V_{NN\pi}\rho V_{NN\pi}^{\dagger}
 +V_{NN}\rho V_{NN\pi}^{\dagger}
 +V_{NN\pi}\rho V_{NN}^{\dagger}.
 \label{eq:coherent-density}
\end{equation}
Tracing the two sectors separately before the cross terms are formed removes
transition interference.  The benchmark must show a nonzero observable error
from that premature trace.

\subsection{Positivity and its limits}

Forward reduced density matrices obey the appropriate positivity tests.
Off-forward GTMD matrices and Wigner distributions are not pointwise positive
objects.  The CP layer cannot be used to clip a negative Wigner
quasidistribution or to replace unresolved coherent amplitudes.

\section{Tagged and differential diagnostics}
\label{sec:tagged}

\subsection{Tagged nucleon and pion channels}

The N1 state permits differential observables in which a spectator nucleon or
pion is retained.  A tagged object stores
\begin{equation}
 \mathfrak t=
 \left(
 \text{tagged species},
 y_s,\bm p_{sT},\lambda_s,c_s,
 \text{active sector},
 \text{FSI status},
 \text{acceptance status},
 \text{inclusive-sum rule}
 \right).
 \label{eq:tagged-id}
\end{equation}
Detector acceptance is a process-layer map and is not included in the nuclear
amplitude.

\subsection{Tagged-to-inclusive closure}

The inclusive N1 parent is recovered by summing all physically unresolved
tagged channels:
\begin{equation}
 W_{a/D}^{N1,\mathrm{incl}}
 =
 \sum_{s=N,\pi}
 \sum_{\lambda_s,c_s}
 \int[\dd\Gamma_s]
 W_{a/D}^{N1,\mathrm{tag}(s)}
 +W_{a/D}^{N1,\mathrm{trans,unresolved}}.
 \label{eq:tagged-inclusive}
\end{equation}
The transition contribution is included according to the actual measurement
resolution.  It cannot be discarded simply because it has no diagonal
probability interpretation.

\subsection{Pole and regular final-state terms}

For nucleon tagging, the analytic on-shell pole remains a useful validation
oracle \cite{CosynWeiss2020,CosynWeiss2026II}.  Regular tagged final-state
interactions may alter the finite-recoil distribution but cannot change the
pole residue.  N1 records tagged FSI as unavailable; it does not absorb such
effects into the nuclear wave function.

\subsection{Tensor tagging}

Tensor-polarized tagged combinations can isolate D-wave and transition
amplitudes.  The benchmark set includes configurations in which the
\(NN\leftrightarrow NN\pi\) interference changes sign while the inclusive
normalization remains stable.  The tagged observable is therefore a
differential test of amplitude content, not merely a normalization check.

\section{Nuclear tensor network and assumption compiler}
\label{sec:tn}

\subsection{Two-sector network topology}

The state network is
\begin{equation}
 \cN_{N1}^{\mathrm{state}}
 =\cN_{NN}\oplus\cN_{NN\pi}.
 \label{eq:n1-state-network}
\end{equation}
The \(NN\) branch couples the correlated proton and neutron microscopic states
to \(S=1\), relative \(L=0,2\), total \(J=1\), and \(I=0\).  The \(NN\pi\)
branch contains the charge-complete isospin tree of
\cref{eq:isospin-zero-state}, the Pauli-allowed nucleon pair, the pion orbital
edge, and the total \(J^P=1^+\) projector.

\subsection{Operator networks}

Separate operator networks represent:
\begin{itemize}
\item the block Hamiltonian;
\item nucleon-active and pion-active partonic operators;
\item transition operators;
\item current and charge operators;
\item coherent diffractive amplitudes and propagation;
\item overlap subtraction maps;
\item partial traces.
\end{itemize}
A physical edge is never reused to represent a provenance relation or Wilson
path merely because the graph topology looks similar.

\subsection{Full-bond and reduced-bond convergence}

At full allowed bond support, the TTN must reproduce direct state and operator
contractions.  Reduced-bond studies report errors in
\begin{equation}
 Z_{NN\pi},
 \quad
 \langle y_{\pi}\rangle,
 \quad
 W^{\pi},
 \quad
 W^{\mathrm{trans}},
 \quad
 b_1,
 \quad
 J_{(2)}^{\mu},
 \quad
 \delta_{\mathrm{ang}},
 \quad
 \delta W_{\mathrm{tensor}}^{\mathrm{coh}}.
 \label{eq:n1-bond-observables}
\end{equation}
A low-rank state that preserves the mass but removes transition interference
or tensor response is not converged.

\subsection{Assumption branches}

The compiler supports mutually exclusive N1 plans such as:
\begin{longtable}{L{0.30\textwidth}L{0.62\textwidth}}
\toprule
Plan & Content\\
\midrule
\endhead
\texttt{N1-PLAN-A} & AV18 N0 state, explicit \(NN\pi\), matched pion subtraction,
consistent current, coherent pilot disabled.\\
\texttt{N1-PLAN-B} & Norfolk N0 state, explicit \(NN\pi\), matched pion subtraction,
consistent current, coherent pilot disabled.\\
\texttt{N1-PLAN-C} & AV18 N0 state, explicit \(NN\pi\), coherent small-\(x\) pilot enabled
with explicit parton--nuclear overlap subtraction.\\
\texttt{N1-INDUCED-REFERENCE} & N0 state with the \(NN\pi\) sector eliminated and
transformed induced operators; read-only comparison.\\
\texttt{N1-ANALYTIC-ORACLE} & Weak-coupling two-sector state and analytic coherent
kernel; validation only.\\
\bottomrule
\end{longtable}
These theories may be compared but never summed.  The coherent pilot is a
branch choice, not an automatic correction.

\section{Global ledgers, covariance, and provenance}
\label{sec:ledgers}

\subsection{State and charge ledgers}

N1 must close separately:
\begin{align}
 Z_{NN}+Z_{NN\pi}&=1,
 \label{eq:ledger-prob}\\
 B_{NN}+B_{NN\pi}&=2,
 \label{eq:ledger-baryon}\\
 Q_{NN}+Q_{NN\pi}&=1,
 \label{eq:ledger-charge}\\
 \langle y_N\rangle+\langle y_\pi\rangle&=1,
 \label{eq:ledger-momentum}\\
 I_{\mathrm{tot}}&=0,
 \qquad
 J^P=1^+.
 \label{eq:ledger-quantum}
\end{align}
The sector contributions are resolution dependent, but the complete ledgers
must close without after-the-fact normalization.

\subsection{Partonic and tensor ledgers}

The common parent must close baryon number, charge, momentum, and the selected
current moments across nucleon-active, pion-active, transition, induced, and
coherent contributions.  The tensor ledger records
\begin{equation}
 \delta_Tf^{N1}
 =\delta_Tf^{NN}
 +\delta_Tf^{NN\pi}
 +\delta_Tf^{\mathrm{trans}}
 +\delta_Tf^{\mathrm{coh}},
 \label{eq:tensor-ledger}
\end{equation}
with no independent normalization attached to the total.

\subsection{Uncertainty decomposition}

The N1 uncertainty record keeps separate:
\begin{enumerate}
\item microscopic proton/neutron member uncertainty;
\item NN interaction and wave-function uncertainty;
\item NNpi interaction and sector-truncation uncertainty;
\item pion partonic-operator uncertainty;
\item internal--exchange separation and overlap-subtraction uncertainty;
\item current and induced-operator truncation;
\item coherent diffractive-input and propagation uncertainty;
\item parton--nuclear soft-overlap uncertainty;
\item TTN bond, quadrature, and discretization errors;
\item future LF-to-QCD matching, evolution, and process uncertainties.
\end{enumerate}
An envelope of alternative potentials is a controlled model family, not
automatically a statistical confidence interval.

\subsection{Provenance two-complex}

The N1 \ProvComplex contains:
\begin{itemize}
\item 0-cells for states, operators, schemes, currents, and results;
\item 1-cells for sector transitions, projections, matching, tracing, and
composition;
\item 2-cells for internal/exchange pion subtraction, explicit/induced
replacement, count-once cuts, current completion, and parton--nuclear overlap
subtraction.
\end{itemize}
Nontrivial homology is an audit signal for an unresolved cycle, missing
subtraction, or genuine alternative mechanism.  It is not a physical
amplitude.

\subsection{Closure vector}

A compact N1 closure vector is
\begin{equation}
 \bm\epsilon_{N1}
 =
 \left(
 \epsilon_{\mathrm{norm}},
 \epsilon_{P^+},
 \epsilon_I,
 \epsilon_J,
 \epsilon_{\mathrm{recoil}},
 \epsilon_{\mathrm{Herm}},
 \epsilon_{\mathrm{current}},
 \epsilon_{\mathrm{ang}},
 \epsilon_{\mathrm{marginal}},
 \epsilon_{b_1},
 \epsilon_{\pi\mathrm{sub}},
 \epsilon_{\mathrm{coh}},
 \epsilon_{\mathrm{CP}},
 \epsilon_{\mathrm{TN}}
 \right).
 \label{eq:n1-closure-vector}
\end{equation}
No component may be hidden inside another uncertainty category.

\section{Computational realization and typed interfaces}
\label{sec:software}

\subsection{Core objects}

\begin{longtable}{L{0.34\textwidth}L{0.58\textwidth}}
\toprule
Object & Responsibility\\
\midrule
\endhead
\texttt{ThreeBodyNuclearCoordinates} & Positive fractions, intrinsic momenta,
Jacobi tree, measure, support, endpoint policy, and normalization.\\
\ThreeBodyRecoil & Active-constituent and spectator recoil at nonzero transfer,
with unit-Jacobian and physical-momentum checks.\\
\TransitionRecoil & Source/target sector kinematics for
\(NN\leftrightarrow NN\pi\) operators.\\
\PionSector & Charge-complete isospin basis, spin/orbital couplings, parity,
Pauli constraints, regulator, and resolution.\\
\texttt{NNPiLightFrontAmplitude} & Spin-, isospin-, charge-, orbital-, and
momentum-resolved three-body amplitudes.\\
\texttt{PionNucleonVertex} & Hermitian sector-changing interaction, momentum
routing, spin/isospin couplers, regulator, and counterterms.\\
\PionOperator & Pion-active quark, antiquark, and gluon operator blocks with
source and matching status.\\
\TransitionOperator & Partonic or current operator between \(NN\) and
\(NN\pi\), including recoil and adjoint rules.\\
\PionSubtraction & Internal/exchange separation projector, overlap term,
remainder, and two-cell identity.\\
\TwoBodyCurrent & Nucleon, pion-in-flight, transition, contact, induced, and
counterterm current pieces.\\
\ContinuityLedger & Signed continuity residual by current and Hamiltonian
component.\\
\texttt{DiffractiveHelicityAmplitude} & Elementary helicity amplitude,
intermediate-state identity, phase, normalization, and source.\\
\CoherentKernel & Longitudinal ordering, propagation phase, scattering order,
helicity, color, and attenuation status.\\
\OverlapMap & Shared region between partonic Wilson and nuclear coherent
propagation.\\
\CoherentBundle & Full matrix-level coherent amplitude before spin projection
or trace.\\
\StateBundle & Normalized state, Hamiltonian, currents, operators, member,
ledgers, solver, and convergence information.\\
\ProvComplex & Typed alternatives, replacements, subtractions, count-once
relations, and observable ancestry.\\
\bottomrule
\end{longtable}

\subsection{Typed data flow}

The principal data flow is
\begin{equation}
 \begin{aligned}
 \texttt{N0StateBundle}
 &\longrightarrow
 \texttt{N1Hamiltonian}
 \longrightarrow
 \StateBundle
 \\
 &\longrightarrow
 \texttt{MatchedN1Operator}
 \longrightarrow
 \texttt{N1DeuteronGTMDParent}
 \\
 &\longrightarrow
 \texttt{ReductionMap}
 \longrightarrow
 \texttt{ValidationObservable}.
 \end{aligned}
 \label{eq:n1-data-flow}
\end{equation}
Every arrow validates coordinates, sectors, operator type, path, Wilson order,
current family, microscopic member, resolution, and provenance.

\subsection{Cache identity}

A parent cache key contains
\begin{equation}
 \mathrm{key}
 =\mathrm{hash}
 \left(
 \mathfrak m_N,
 \mathfrak m_D,
 r_N,
 \mathfrak o_{a/D}^{N1},
 \mathfrak R_{N1},
 x_D,\kT,\DT,
 \text{solver},
 \text{code version}
 \right).
 \label{eq:n1-cache-key}
\end{equation}
Changing the wave function, pion separator, current, path, coherent branch,
recoil, or member invalidates the cache.

\section{Formal requirements and property-based tests}
\label{sec:requirements}

\subsection{Stable requirement identifiers}

The following requirement families are stable machine-readable contracts.
\begin{longtable}{L{0.16\textwidth}L{0.76\textwidth}}
\toprule
ID family & Requirement\\
\midrule
\endhead
\texttt{V13.SCOPE.*} & N1 remains a validation-only \(NN\oplus NN\pi\) theory
with all downstream readiness gates explicit.\\
\texttt{V13.STATE.*} & One normalized state, exact sector identities, one
charge and momentum ledger, and a Hermitian block Hamiltonian.\\
\texttt{V13.ISO.*} & Charge-complete \(I=0\) coupling, exact CG phases, and
separate CSB branches.\\
\texttt{V13.PARITY.*} & Pion intrinsic parity, odd total orbital parity, Pauli
constraints, and total \(J^P=1^+\).\\
\texttt{V13.KIN.*} & Three-body support, Jacobi maps, diagonal recoil,
transition recoil, unit Jacobians, and endpoint policy.\\
\texttt{V13.OP.*} & Complete sector matrix, pion-active, nucleon-active,
transition, contact, and induced operator identities.\\
\texttt{V13.PARENT.*} & Amplitude-level sector sum, Hermiticity, link identity,
spin-1 reconstruction, and common-member ancestry.\\
\texttt{V13.PION.*} & Direct pion-active reductions, positive-\(x\) species
identity, transition interference, and moment closure.\\
\texttt{V13.SUB.*} & Internal/exchange separator, overlap subtraction,
explicit/induced equivalence, remainder, and no-double-counting two-cells.\\
\texttt{V13.CURR.*} & Hamiltonian-consistent current family, continuity ledger,
angular condition, GTMD/current closure, and ablation residuals.\\
\texttt{V13.COH.*} & Helicity-resolved diffractive amplitudes, propagation,
scalar/vector/tensor separation, zero limits, and parton--nuclear subtraction.\\
\texttt{V13.CP.*} & Amplitude embedding before trace, complete positivity,
trace closure, and premature-trace failure.\\
\texttt{V13.TAG.*} & Tagged identities, inclusive recovery, pole oracle,
transition resolution, and acceptance separation.\\
\texttt{V13.TN.*} & Two-sector TTN, full-bond equality, observable-sensitive
bond convergence, and symmetry-block retention.\\
\texttt{V13.PROV.*} & Typed alternatives, replacements, subtractions,
count-once relations, deterministic normalization, and ancestry queries.\\
\texttt{V13.CONV.*} & Multi-axis basis, sector, endpoint, interaction, current,
quadrature, coherent-order, and bond convergence.\\
\bottomrule
\end{longtable}

\subsection{Kinematic and state tests}

Mandatory tests include:
\begin{enumerate}
\item positive three-body support and exact fraction closure;
\item intrinsic transverse closure in both momentum fibers;
\item unit Jacobian for diagonal and transition maps;
\item full physical transfer to the active constituent and unchanged
spectator momenta;
\item charge-complete \(I=0\) reconstruction;
\item exact total charge, baryon number, parity, \(J^z\), and Pauli closure;
\item Hamiltonian Hermiticity and generated sector adjoint;
\item exact/Krylov/full-bond agreement;
\item state normalization and plus-momentum closure;
\item state tracking across the N1 resolution tower.
\end{enumerate}

\subsection{Operator and reduction tests}

The operator suite tests:
\begin{enumerate}
\item source/target sector and active-constituent compatibility;
\item diagonal nucleon-active and pion-active blocks;
\item transition-operator Hermiticity and recoil;
\item complete target/parton helicity matrices;
\item spin-1 projector reconstruction and tensor sign adapter;
\item GTMD-to-TMD/GPD/PDF/current closure;
\item pion-active number and momentum moments;
\item transition-interference zero limits;
\item future/past and gluon \(f/d\) identity preservation;
\item absence of nuclear coherent phases from the partonic path layer.
\end{enumerate}

\subsection{Current and coherence tests}

The current and coherent suite tests:
\begin{enumerate}
\item continuity closure with all current pieces;
\item signed residuals after each required current ablation;
\item charge, magnetic, quadrupole, and angular-condition closure;
\item direct-current versus integrated-parent equality;
\item zero coherent correction when either elementary amplitude vanishes;
\item longitudinal-order reversal;
\item distinct scalar, vector, and tensor coherent projections;
\item failure of a copied unpolarized shadowing ratio;
\item one correct parton--nuclear overlap subtraction;
\item signed defects for missing or duplicate subtraction.
\end{enumerate}

\subsection{Provenance and CP tests}

The provenance and reduction suite tests:
\begin{enumerate}
\item rejection of internal plus exchange pion without overlap subtraction;
\item rejection of explicit \(NN\pi\) plus its induced replacement;
\item visible Feshbach remainder and norm kernel;
\item count-once equivalent cut and overlap paths;
\item additive retention of physically distinct channels;
\item exact equality of amplitude-derived and Kraus CP maps;
\item nonzero error from tracing before transition/coherent amplitudes are
combined;
\item deterministic path normalization and serialization;
\item no route to production, matching, evolution, process, or inference
roots.
\end{enumerate}

\subsection{Required negative-test scale}

A production-quality N1 implementation should contain at least 280 stable
negative-test identities.  The catalogue must cover:
\begin{itemize}
\item incomplete charge or isospin bases;
\item wrong pion parity, Pauli phase, or angular coupling;
\item reused two-body coordinates in the three-body sector;
\item broken active/spectator recoil or transition Jacobian;
\item diagonal/transition sector aliases;
\item copied proton/neutron or quark/antiquark parents;
\item missing pion-active or transition operator identity;
\item double-counted internal/exchange/induced pion mechanisms;
\item current omissions and independently normalized current components;
\item copied tensor shadowing from an unpolarized ratio;
\item premature partial trace;
\item hidden coherent or Wilson path substitution;
\item reduced-bond loss hidden by an energy-only test;
\item unsupported \(\Delta\Delta\), compact, shadowing, matching, evolution,
process, inference, or production promotion.
\end{itemize}

\section{Analytic and finite-dimensional benchmark suite}
\label{sec:benchmarks}

\subsection{N1-A: charge-complete isospin state}

Construct the three charge channels in \cref{eq:charge-basis}, couple them to
\(I=0\), and verify normalization, phase convention, total charge, and total
isospin.  Removing any one charge channel must produce a nonzero isospin
residual.

\subsection{N1-B: parity and Pauli basis}

Build the minimal \(J^P=1^+\) \(NN\pi\) basis with odd
\(\ell_{NN}+\ell_\pi\).  Verify \cref{eq:pauli-nnpi}.  Injecting an even total
orbital parity or a Pauli-forbidden pair state must fail.

\subsection{N1-C: three-body recoil}

Use an analytic three-body Gaussian amplitude and verify intrinsic closure,
physical spectator preservation, full active transfer, reversal, and unit
Jacobian for active nucleon and active pion choices.  Independently compare
against a direct physical-momentum construction.

\subsection{N1-D: two-sector Hamiltonian}

Use a finite Hermitian \(NN\oplus NN\pi\) model.  Compare exact, matrix-free,
Krylov, and full-bond solutions.  Demonstrate normalization and parameter flow
at several resolutions.

\subsection{N1-E: pion-active common-parent closure}

Choose a regulated spin-zero pion parent and verify direct and sequential
TMD/GPD/PDF/current routes at forward and nonzero transfer.  The pion oracle
is validation only and cannot be promoted to a physical TMD.

\subsection{N1-F: transition interference}

Construct a state for which the diagonal pion contribution is small but the
\(NN\leftrightarrow NN\pi\) interference is nonzero.  Verify its phase,
Hermiticity, tensor projection, and exact zero when either sector amplitude or
the transition operator is disabled.

\subsection{N1-G: explicit versus induced pion}

Eliminate the \(NN\pi\) sector and compare the full matrix element with the
induced Hamiltonian and transformed operator.  The equivalence residual must
close while a nonzero remainder and norm kernel remain visible.  Using
\(POP\) alone must fail.

\subsection{N1-H: internal--exchange subtraction}

Construct overlapping internal and exchange pion regions.  Verify
\cref{eq:pion-matched-sum} at several separation scales.  Missing subtraction
and duplicate subtraction must give equal-and-opposite signed defects.

\subsection{N1-I: current continuity}

Build a finite current with nucleon, pion-in-flight, transition, contact,
induced, and counterterm pieces.  Verify \cref{eq:n1-continuity}.  Remove each
required nonzero term and record the signed residual.

\subsection{N1-J: spin-1 current and moments}

Generate spin-1 current helicity amplitudes from chosen form factors and a
consistent two-sector state.  Verify charge, magnetic, quadrupole, angular
condition, and integrated GTMD/current closure.

\subsection{N1-K: coherent double scattering}

Use a two-state diffractive kernel with known propagation phase.  Test zero
limits, ordering reversal, scalar/vector/tensor differences, and failure of a
copied unpolarized ratio.

\subsection{N1-L: parton--nuclear overlap}

Construct a shared soft region between a partonic Wilson term and the coherent
nuclear pilot.  Verify one subtraction, signed missing/duplicate residuals,
and strict preservation of the distinct path identities.

\subsection{N1-M: CP reduction}

Construct an amplitude embedding with resolved \(NN\) and \(NN\pi\) coherence.
Compare direct partial trace with the Kraus representation.  Then trace the
sectors separately before interference is formed and demonstrate the wrong
result.

\subsection{N1-N: tagged-inclusive closure}

Implement tagged nucleon and pion channels, sum them to the inclusive parent,
and verify the analytic pole oracle.  A regular FSI test term may modify the
finite tagged distribution but not the pole residue.

\subsection{N1-O: tensor-network sensitivity}

Compare full-bond and reduced-bond networks.  Require at least one reduced bond
to reproduce the mass within a small error while losing a substantial
pion-transition, \(b_1\), current, or coherent tensor signal.

\subsection{N1-P: assumption and provenance plans}

Compile all supported N1 plans.  Verify that explicit and induced pion
branches, alternative nuclear wave functions, and coherent-on/off branches
are mutually exclusive where required and that each result has a unique
content hash and ancestry.

\section{Staged N1 implementation program}
\label{sec:stages}

\subsection{N1.0: types and kinematics}

Implement the charge-complete sector, three-body coordinates, recoil maps,
transition map, parity and Pauli projectors, state and operator identities,
and analytic recoil benchmarks.

\subsection{N1.1: two-sector state}

Implement the Hermitian Hamiltonian, small resolution tower, exact/Krylov/TTN
solvers, state tracking, normalization and momentum ledgers, and the analytic
weak-coupling oracle.

\subsection{N1.2: pion-active and transition operators}

Implement diagonal nucleon-active and pion-active blocks, transition
operators, common-parent helicity matrices, spin-1 projectors, and regulated
reduction closure.

\subsection{N1.3: pion matching and currents}

Implement the internal/exchange separator, overlap subtraction, Feshbach
comparison, Hamiltonian-consistent current family, continuity ledger, and
current/GTMD moment closure.

\subsection{N1.4: coherent and CP pilots}

Implement helicity-resolved diffractive amplitudes, propagation, tensor
projections, parton--nuclear overlap subtraction, amplitude-derived CP maps,
and premature-trace failure.

\subsection{N1.5: convergence and release gates}

Complete multi-axis convergence, tagged diagnostics, tensor-network bond
studies, deterministic manifests, provenance normalization, source-hash audit,
and downstream fail-closed tests.

\section{Acceptance criteria and scientific status}
\label{sec:acceptance}

The N1 formalism is internally accepted only when:
\begin{enumerate}
\item the \(NN\pi\) sector has complete charge, isospin, spin, parity, orbital,
and momentum identity;
\item the two-sector Hamiltonian is Hermitian and solved consistently across a
resolution tower;
\item one normalized state closes probability, charge, baryon, isospin, and
plus momentum;
\item diagonal and transition recoil maps close and have declared Jacobians;
\item pion-active, nucleon-active, and transition operator blocks are complete
and compose only with compatible sectors;
\item the complete N1 parent satisfies Hermiticity and reconstructs all target
and parton helicity matrices;
\item GTMD, TMD, GPD, PDF, current, and tensor reductions close from the same
parent;
\item the internal/exchange pion subtraction is executable and stable under
separator variation within truncation errors;
\item explicit and induced pion descriptions are related by transformed
operators and a visible remainder;
\item the Hamiltonian-consistent current satisfies the finite-basis continuity
benchmark and the spin-1 angular condition;
\item pion-active and transition contributions to \(b_1\) are predicted by the
shared state and not independently normalized;
\item the coherent pilot is helicity resolved, has exact zero and ordering
limits, and does not copy an unpolarized ratio;
\item the partonic Wilson and nuclear coherent sectors remain distinct and any
shared region is subtracted once;
\item the CP map is derived only after coherent amplitudes are combined;
\item tagged integration closes to the inclusive parent;
\item exact, Krylov, full-bond, and reduced-bond routes satisfy their declared
convergence obligations;
\item the provenance complex rejects every explicit/induced, internal/exchange,
cut, current, and coherent double count;
\item all unavailable sectors and downstream matching, evolution, process, and
inference gates remain fail closed;
\item no accepted production artifact or registry route is modified;
\item at least one tensor, pion-active, transition, current, and coherent
observable remains withheld from calibration.
\end{enumerate}

\begin{decision}[N1 predictive status]
Passing the N1 gates establishes a regulated, spin-resolved
\(NN\oplus NN\pi\) microscopic validation parent.  It does not establish a
physical deuteron TMD or physical shadowing prediction.  The next dynamical
layer must add supported \(\Delta\Delta\), compact six-quark/hidden-color, or
higher coherent sectors only through complete amplitudes and matched
operators.
\end{decision}

\section{Boundary to N2 and later matching}
\label{sec:boundary}

A successful N1 package exports:
\begin{itemize}
\item the normalized N1 state bundle and resolution trajectory;
\item complete \(NN\) and \(NN\pi\) amplitudes;
\item pion-active, nucleon-active, transition, and current operators;
\item the matched internal/exchange pion ledger;
\item order-resolved quark, antiquark, and gluon deuteron parents;
\item the coherent pilot and parton--nuclear overlap ledger;
\item exact/full-bond and reduced-bond convergence manifests;
\item tagged and tensor validation observables;
\item the complete provenance two-complex.
\end{itemize}

The likely next nuclear package is
\begin{quote}
\textbf{N2: explicit \(\Delta\Delta\) and compact six-quark/hidden-color
sectors, transition currents, normalized non-nucleonic interference, and an
upgraded coherent amplitude.}
\end{quote}
A narrower N1 completion package is preferred if current continuity, pion
subtraction, or coherent-overlap residuals remain parametrically larger than
the controlled nuclear-sector truncation.

\section{Conclusion}

The \(NN\pi\) sector is the first nuclear extension for which the distinction
between a mechanism label and a genuine amplitude becomes unavoidable.  A
pion distribution, a meson-exchange current, a transition interference, and a
coherent shadowing correction are not four independent curves.  They are
sector-diagonal, sector-changing, current, and multiple-scattering matrix
elements of one regulated nuclear theory.  Volume XIII makes that statement
computationally precise.  It supplies the charge-complete three-body state,
exact recoil geometry, matched operator matrix, internal--exchange
subtraction, consistent current family, helicity-resolved coherent pilot,
derived CP map, tensor-network convergence criteria, and provenance relations
needed for C16/N1.  The result is the first formal nuclear layer beyond
impulse approximation that preserves the microscopic quark--antiquark--gluon
parent without rebuilding a patchwork at the deuteron level.

\appendix

\section{Derivation of the three-body spectator-preserving recoil map}
\label{app:recoil}

Let the average intrinsic momenta satisfy
\(
\sum_i\bm k_{iT}=0
\).
For active constituent \(j\), impose
\begin{equation}
 p_{iT}^{\mathrm{out}}=p_{iT}^{\mathrm{in}}
 \qquad(i\neq j),
 \label{eq:app-spectator}
\end{equation}
with
\begin{equation}
 p_{iT}^{\mathrm{in/out}}
 =y_iP_{i/f,T}+\bm\kappa_{iT}^{\mathrm{in/out}}.
 \label{eq:app-physical}
\end{equation}
Using \(P_{iT}=-\DT/2\) and \(P_{fT}=+\DT/2\) gives
\begin{equation}
 \bm\kappa_{iT}^{\mathrm{out}}
 -\bm\kappa_{iT}^{\mathrm{in}}
 =-y_i\DT
 \qquad(i\neq j).
 \label{eq:app-spectator-shift}
\end{equation}
Intrinsic closure then gives
\begin{equation}
 \bm\kappa_{jT}^{\mathrm{out}}
 -\bm\kappa_{jT}^{\mathrm{in}}
 =(1-y_j)\DT.
 \label{eq:app-active-shift}
\end{equation}
The symmetric choice around the average variables yields
\cref{eq:active-threebody-recoil,eq:spectator-threebody-recoil}.  Substitution
into \cref{eq:app-physical} gives
\(
 p_{jT}^{\mathrm{out}}-p_{jT}^{\mathrm{in}}=\DT
\)
and unchanged spectators.

\section{Isospin and charge adapters}
\label{app:isospin}

Use
\begin{equation}
 |1,+1\rangle_{NN}=|pp\rangle,
 \qquad
 |1,0\rangle_{NN}=\frac{|pn\rangle+|np\rangle}{\sqrt2},
 \qquad
 |1,-1\rangle_{NN}=|nn\rangle.
 \label{eq:triplet-nn}
\end{equation}
The pion triplet is
\begin{equation}
 |1,+1\rangle_{\pi}=|\pi^+\rangle,
 \quad
 |1,0\rangle_{\pi}=|\pi^0\rangle,
 \quad
 |1,-1\rangle_{\pi}=|\pi^-\rangle.
 \label{eq:pion-triplet}
\end{equation}
Then \cref{eq:isospin-zero-state} follows from the standard
\(1\otimes1\to0\) coefficient.  A software adapter stores these phases and
verifies unitarity under conversion between coupled and charge bases.

\section{Finite two-sector continuity benchmark}
\label{app:continuity}

Let
\begin{equation}
 H=\begin{pmatrix}H_P&V^{\dagger}\\V&H_Q\end{pmatrix},
 \qquad
 \rho=\begin{pmatrix}\rho_P&\rho_{PQ}\\\rho_{QP}&\rho_Q\end{pmatrix}.
 \label{eq:two-sector-h-rho}
\end{equation}
The continuity identity decomposes into diagonal and transition blocks.  For
example,
\begin{equation}
 q_{\mu}J_{PQ}^{\mu}
 =H_P\rho_{PQ}+V^{\dagger}\rho_Q
 -\rho_PV^{\dagger}-\rho_{PQ}H_Q.
 \label{eq:transition-continuity}
\end{equation}
A transition current is therefore required whenever the sector-changing
interaction and charge operator do not commute.  A one-body current alone
cannot satisfy the complete identity.

\section{Machine-readable schemas}
\label{app:schemas}

\subsection{N1 state record}

A state record contains:
\begin{verbatim}
{
  "state_id": "...",
  "resolution_id": "...",
  "assumption_plan": "...",
  "sectors": ["NN", "NNPI"],
  "nn_member_id": "...",
  "nnpi_amplitude_id": "...",
  "charge_isospin_basis_id": "...",
  "hamiltonian_id": "...",
  "renormalization_trajectory_id": "...",
  "solver_id": "...",
  "tensor_network_manifest_id": "...",
  "normalization_ledger_id": "...",
  "readiness": ["N1_STATE_VALIDATION_ONLY"]
}
\end{verbatim}

\subsection{Pion subtraction record}

\begin{verbatim}
{
  "subtraction_id": "...",
  "internal_pion_source": "...",
  "exchange_pion_source": "...",
  "separator_id": "...",
  "overlap_operator_id": "...",
  "matched_result_id": "...",
  "separation_scale": "...",
  "remainder_id": "...",
  "two_cell_id": "..."
}
\end{verbatim}

\subsection{Coherent amplitude record}

\begin{verbatim}
{
  "coherent_id": "...",
  "parton_species": "...",
  "operator_projection": "...",
  "initial_target_helicity": "...",
  "final_target_helicity": "...",
  "intermediate_state": "...",
  "elementary_amplitudes": ["...", "..."],
  "propagation_kernel_id": "...",
  "longitudinal_order": "...",
  "scattering_order": 2,
  "parton_nuclear_overlap_id": "...",
  "factorization_status": "VALIDATION_ONLY"
}
\end{verbatim}

\section{Compact N1 acceptance checklist}

\begin{enumerate}
\item Complete \(NN\pi\) charge/isospin basis?
\item Correct pion parity and Pauli-allowed orbital channels?
\item Three-body and transition recoil maps validated?
\item One normalized two-sector state and momentum ledger?
\item Hermitian Hamiltonian and generated transition adjoint?
\item Complete pion-active, nucleon-active, and transition operators?
\item Common helicity parent and commuting reductions?
\item Explicit internal/exchange pion overlap subtraction?
\item Explicit/induced equivalence with transformed operators and remainder?
\item Hamiltonian-consistent current and continuity ledger?
\item Spin-1 angular condition and current/GTMD moments?
\item Helicity-resolved coherent pilot with exact zero limits?
\item Partonic/nuclear soft overlap subtracted once?
\item CP map derived only after amplitudes are combined?
\item Tagged-inclusive and pole-oracle closure?
\item Full-bond equality and observable-sensitive bond convergence?
\item Deterministic provenance two-complex and ancestry?
\item All unavailable sectors and downstream gates fail closed?
\item Production and authoritative artifacts unchanged?
\item At least one nontrivial pion/tensor/current/coherent holdout preserved?
\end{enumerate}

\small
\begin{thebibliography}{99}

\bibitem{Dirac1949}
P.~A.~M.~Dirac,
``Forms of relativistic dynamics,''
Rev.\ Mod.\ Phys.\ \textbf{21}, 392 (1949).

\bibitem{KeisterPolyzou1991}
B.~D.~Keister and W.~N.~Polyzou,
``Relativistic Hamiltonian dynamics in nuclear and particle physics,''
Adv.\ Nucl.\ Phys.\ \textbf{20}, 225 (1991).

\bibitem{Carbonell1998}
J.~Carbonell, B.~Desplanques, V.~A.~Karmanov, and J.-F.~Mathiot,
``Explicitly covariant light-front dynamics and relativistic few-body systems,''
Phys.\ Rept.\ \textbf{300}, 215 (1998), arXiv:nucl-th/9804029.

\bibitem{Wiringa1995}
R.~B.~Wiringa, V.~G.~J.~Stoks, and R.~Schiavilla,
``Accurate nucleon-nucleon potential with charge-independence breaking,''
Phys.\ Rev.\ C \textbf{51}, 38 (1995), arXiv:nucl-th/9408016.

\bibitem{Machleidt2001}
R.~Machleidt,
``The high-precision, charge-dependent Bonn nucleon-nucleon potential,''
Phys.\ Rev.\ C \textbf{63}, 024001 (2001), arXiv:nucl-th/0006014.

\bibitem{Piarulli2016}
M.~Piarulli et al.,
``Local chiral potentials and the structure of light nuclei,''
Phys.\ Rev.\ C \textbf{94}, 054007 (2016), arXiv:1606.06335.

\bibitem{Krebs2020}
H.~Krebs,
``Nuclear currents in chiral effective field theory,''
Eur.\ Phys.\ J.\ A \textbf{56}, 234 (2020), arXiv:2008.00974.

\bibitem{Baroni2015}
A.~Baroni, L.~Girlanda, S.~Pastore, R.~Schiavilla, and M.~Viviani,
``Nuclear axial currents in chiral effective field theory,''
Phys.\ Rev.\ C \textbf{93}, 015501 (2016), arXiv:1509.07039.

\bibitem{Schiavilla2018}
R.~Schiavilla et al.,
``Local chiral interactions and magnetic structure of few-nucleon systems,''
Phys.\ Rev.\ C \textbf{99}, 034005 (2019), arXiv:1809.10180.

\bibitem{HeZahed2024}
F.~He and I.~Zahed,
``Deuteron gravitational form factors: Exchange currents,''
arXiv:2401.09318.

\bibitem{Sullivan1972}
J.~D.~Sullivan,
``One-pion exchange and deep-inelastic electron-nucleon scattering,''
Phys.\ Rev.\ D \textbf{5}, 1732 (1972).

\bibitem{Miller2014}
G.~A.~Miller,
``Pionic and hidden-color, six-quark contributions to the deuteron \(b_1\)
structure function,''
Phys.\ Rev.\ C \textbf{89}, 045203 (2014), arXiv:1311.4561.

\bibitem{CaoEtAl2026}
X.~Cao, S.~Cheng, Y.~Duan, Y.~Li, S.~Xu, and X.~Zhao,
``Non-perturbative flavor asymmetry in the nucleon and deuteron: The
light-front Hamiltonian effective field theory approach,''
arXiv:2601.13567 (2026).

\bibitem{FrankfurtGuzeyStrikman2012}
L.~Frankfurt, V.~Guzey, and M.~Strikman,
``Leading twist nuclear shadowing phenomena in hard processes with nuclei,''
Phys.\ Rept.\ \textbf{512}, 255 (2012), arXiv:1106.2091.

\bibitem{HoodbhoyJaffeManohar1989}
P.~Hoodbhoy, R.~L.~Jaffe, and A.~Manohar,
``Novel effects in deep-inelastic scattering from spin-one hadrons,''
Nucl.\ Phys.\ B \textbf{312}, 571 (1989).

\bibitem{HERMES2005}
A.~Airapetian et al.\ (HERMES Collaboration),
``First measurement of the tensor structure function \(b_1\) of the
deuteron,''
Phys.\ Rev.\ Lett.\ \textbf{95}, 242001 (2005), arXiv:hep-ex/0506018.

\bibitem{CosynDongKumanoSargsian2017}
W.~Cosyn, Y.-B.~Dong, S.~Kumano, and M.~Sargsian,
``Tensor-polarized structure function \(b_1\) in the standard convolution
description of the deuteron,''
Phys.\ Rev.\ D \textbf{95}, 074036 (2017), arXiv:1702.05337.

\bibitem{CosynWeiss2020}
W.~Cosyn and C.~Weiss,
``Polarized electron--deuteron deep-inelastic scattering with spectator
nucleon tagging,''
Phys.\ Rev.\ C \textbf{102}, 065204 (2020), arXiv:2006.03033.

\bibitem{CosynWeiss2026II}
W.~Cosyn and C.~Weiss,
``Semi-inclusive deep-inelastic scattering on a polarized spin-1 target. II.
Deuteron and spectator nucleon tagging,''
arXiv:2603.23700 (2026).

\bibitem{BacchettaMulders2000}
A.~Bacchetta and P.~J.~Mulders,
``Deep inelastic leptoproduction of spin-one hadrons,''
Phys.\ Rev.\ D \textbf{62}, 114004 (2000), arXiv:hep-ph/0007120.

\bibitem{KumanoSong2021}
S.~Kumano and Q.-T.~Song,
``Transverse-momentum-dependent parton distribution functions for spin-1
hadrons,''
Phys.\ Rev.\ D \textbf{103}, 014025 (2021), arXiv:2011.08583.

\bibitem{Boer2016}
D.~Boer, S.~Cotogno, T.~van Daal, P.~J.~Mulders, A.~Signori, and Y.-J.~Zhou,
``Gluon and Wilson-loop TMDs for hadrons of spin \(\leq1\),''
JHEP \textbf{10}, 013 (2016), arXiv:1607.01654.

\bibitem{Collins2011}
J.~C.~Collins,
\emph{Foundations of Perturbative QCD},
Cambridge University Press (2011).

\end{thebibliography}

\end{document}
